\documentclass[11pt]{article}

\usepackage[margin=1in]{geometry}
\usepackage[T1]{fontenc}
\usepackage[utf8]{inputenc}
\usepackage{lmodern}
\usepackage{microtype}
\usepackage{enumitem}
\usepackage{booktabs}
\usepackage{longtable}
\usepackage{array}
\usepackage{xcolor}
\usepackage{hyperref}
\usepackage{listings}
\usepackage{titlesec}
\usepackage{fancyhdr}

\hypersetup{
  colorlinks=true,
  linkcolor=blue!60!black,
  urlcolor=blue!60!black,
  citecolor=blue!60!black,
  pdftitle={Cold Review of ALCI_RCC5 Lean Development, Rounds 19--25},
  pdfauthor={OpenAI GPT-5.5 Pro}
}

\lstdefinelanguage{LeanPseudo}{
  morekeywords={structure,where,def,theorem,Prop,inductive,abbrev,forall,exists,fun,if,then,else,by,match,with,deriving},
  sensitive=true,
  morecomment=[l]{--}
}
\lstset{
  basicstyle=\ttfamily\small,
  breaklines=true,
  columns=fullflexible,
  keepspaces=true,
  frame=single,
  framerule=0.3pt,
  rulecolor=\color{black!25},
  xleftmargin=0.5em,
  xrightmargin=0.5em,
  language=LeanPseudo
}

\setlist[itemize]{topsep=3pt,itemsep=2pt,parsep=1pt,leftmargin=1.35em}
\setlist[enumerate]{topsep=3pt,itemsep=2pt,parsep=1pt,leftmargin=1.55em}
\renewcommand{\arraystretch}{1.16}
\newcolumntype{L}[1]{>{\raggedright\arraybackslash}p{#1}}
\pagestyle{fancy}
\fancyhf{}
\lhead{Cold review: ALCI\_RCC5 Lean rounds 19--25}
\rhead{\thepage}

\title{Cold Review of the \texttt{ALCI\_RCC5} Lean Development\\Rounds 19--25}
\author{Adversarial adequacy review of definitions and theorem statements}
\date{14 July 2026}

\begin{document}
\maketitle

\begin{abstract}
This report reviews the uploaded artifact \texttt{Round19Transport.lean}
against the neutral reference semantics in \texttt{ALCI\_RCC5\_REFERENCE.md}.
It does not re-check kernel proofs.  The target is adequacy: whether the
formal statements mean ALCI\_RCC5 concept satisfiability and whether the
certificate/catalo\-gue interface is non-vacuous and finite enough to support
the advertised reduction.

The bottom line is \textbf{gap, repairable}.  I agree that the defects appear
fixable and local to the modelling/reduction interface rather than signs that
the kernel-checked proofs are invalid.  I do \emph{not} regard all of them as
minor relative to the advertised decidability story: the soundness pipeline is
credible, but the finite-certificate/completeness interface needs substantive
tightening before it can honestly support a decision-procedure claim.  This is
not a second mathematical conjecture beyond bounded-width/F3 completeness; it
is a formalization-level objection that the current Lean statement does not yet
express that assumption as a bounded, finite, computably checkable obligation.
\end{abstract}

\section{Scope and standard of review}

The submitted source says it is the normative artifact and that the prose
companions are commentary.  I used \texttt{ALCI\_RCC5\_REFERENCE.md} as the
yardstick.  Line numbers in this report refer to the file as archived in
\texttt{cold\_review\_round25.zip}.

The review question is not whether Lean accepted the proofs.  It is whether
the Lean definitions and theorem statements match:
\begin{itemize}
  \item the standard RCC5 atomic-network semantics with strong equality;
  \item NNF ALCI syntax with RCC5 atoms as roles and inverse roles absorbed by
        converse;
  \item concept satisfiability as nonempty extension of the input concept;
  \item a finite and computably searchable certificate/catalogue reduction.
\end{itemize}

\section{Verdict}

\textbf{Verdict: gap, repairable.}

The core soundness chain is largely adequate:
\[
  \texttt{Wellformed} + \texttt{SCond} + \texttt{Hintikka}
  \quad\Longrightarrow\quad
  \texttt{certInterp\_rcc5}
  \quad\Longrightarrow\quad
  \texttt{truth\_lemma}
  \quad\Longrightarrow\quad
  \texttt{Satisfiable}.
\]
I found no polarity error in \texttt{sat}, no missing NNF constructor in
\texttt{Concept}, and no apparent RCC5 composition-table error.

The gap is in the reduction interface.  The current formalization proves a
sound machine from generated objects to models.  It does not yet formalize an
honest finite, bounded, decidable search space for those generated objects.
The most important defects are:
\begin{enumerate}
  \item \texttt{Cert}, \texttt{Template}, and \texttt{Catalog} contain arbitrary
        higher-order function fields, so they are not finite syntactic
        certificates as Lean objects.
  \item \texttt{Satisfiable} is monomorphic over the countable occurrence
        carrier \texttt{Occ}, whereas the reference semantics quantifies over
        arbitrary nonempty domains of regions.
  \item \texttt{CompletenessObligation} is an unbounded existence proposition,
        not a finite bounded search theorem.
  \item \texttt{SCat.net\_r3} is over-strong on degenerate template triples and
        can reject semantically equivalent certificates because of unused
        diagonal template values.
\end{enumerate}

These are not proof failures; they are definitional and statement-level gaps.
They are repairable, but they are substantive for the advertised reduction to
decidability.

\paragraph{Relation to bounded-width/F3.}
The finite-certificate finding should not be read as a new open mathematical
assumption in addition to the bounded-width/F3 completeness claim.  The latter
remains the main unproved mathematical step: every satisfiable concept should
have a bounded generated certificate of the advertised kind.  The formalization
problem is that the current \texttt{CompletenessObligation} does not yet state
that step in a decision-grade form.  It quantifies over higher-order catalogue
objects and an unbounded labelling rather than over finite codes bounded by a
computable function of the input concept.  Thus the defect is serious for the
claim that the Lean artifact has isolated exactly one decision-ready
assumption, but it is not evidence of an additional RCC5 obstruction.

\section{Summary table}

The IDs in the following table are referee finding labels only; in particular,
Finding F3 is unrelated to the project's bounded-width/F3 terminology.

{\small
\begin{longtable}{L{0.055\linewidth}L{0.265\linewidth}L{0.405\linewidth}L{0.185\linewidth}}
\toprule
ID & Location & Finding & Severity \\
\midrule
F1 & \texttt{Template}; \texttt{Cert}; \texttt{Catalog}; \texttt{CatOk.F\_reads}; \texttt{Completeness-}\linebreak\texttt{Obligation} & Certificates/catalogues are represented by arbitrary total and higher-order functions, not finite syntax. & Fatal to finite-decision claim; repairable \\
F2 & \texttt{Interp}; \texttt{Satisfiable} & Satisfiability is over subdomains of \texttt{Occ}, not arbitrary reference domains. & Major adequacy gap; repairable \\
F3 & \texttt{Completeness-}\linebreak\texttt{Obligation} & Obligation lacks computable bounds, finite code, closure restriction, and decidability theorem. & Fatal to current reduction claim; repairable \\
F4 & \texttt{SCat.net\_r3} versus \texttt{SCond.tri} & Catalogue check quantifies degenerate triples involving the fresh port itself; semantic condition excludes them. & Moderate over-restriction; repairable \\
F5 & \texttt{certC}; \texttt{certK}; \texttt{sample-}\linebreak\texttt{satisfiable\_ex} & Positive witnesses are mostly all-DR and do not exercise hard positive composition/catalogue cases. & Evidence gap, not fatal \\
F6 & Headline theorem statements & Main soundness statements are mostly honest; some prose overstates hypotheses or finiteness. & Mostly adequate; wording repairs \\
\bottomrule
\end{longtable}
}

\section{Findings}

\subsection{F1. \texttt{Catalog} and \texttt{Cert} are not finite certificates}

\textbf{Locations.}
\begin{itemize}
  \item \texttt{Template.net}: lines 224--228.
  \item \texttt{Cert.coreNet} and \texttt{Cert.f}: lines 234--241.
  \item \texttt{Catalog.coreNet} and \texttt{Catalog.F}: lines 3311--3317.
  \item \texttt{CatOk.F\_reads}: lines 3505--3514.
  \item \texttt{CompletenessObligation}: lines 4459--4476.
\end{itemize}

The artifact repeatedly uses finite-certificate language, but the Lean
structures contain fields of the following shape:

\begin{lstlisting}
structure Template where
  nSlots : Nat
  nFresh : Nat
  net : TPort -> TPort -> Atom

structure Cert where
  nCore : Nat
  coreNet : Nat -> Nat -> Atom
  templates : List Template
  steps : List Step
  f : Nat -> (Nat -> Atom) -> (Nat -> Atom) -> Nat -> Atom

structure Catalog where
  nCore : Nat
  coreNet : Nat -> Nat -> Atom
  templates : List Template
  root : Nat
  attaches : List Attach
  F : Nat -> (Nat -> Atom) -> (Nat -> Atom) -> Nat -> Atom
\end{lstlisting}

These fields are arbitrary functions.  The field \texttt{F} is higher-order in
row functions.  The later predicates constrain finite slices of their behavior,
but the witnesses themselves are not finite codes.  In particular,
\texttt{CatOk.F\_reads} is an extensionality discipline over arbitrary row
functions, not a finite lookup table.

\paragraph{Why this matters.}
A Lean proof of \texttt{CompletenessObligation} could choose a function
\texttt{F} that hides an infinite or noncomputable construction.  Even if all
queries made by \texttt{SCat} and \texttt{Hintikka} are finite for a fixed
object, the object being existentially quantified is not itself a finite
certificate.  Thus the current statement does not by itself support finite
search.

\paragraph{Concrete symptom.}
Two catalogues may agree on every value ever inspected by a given plan and
\texttt{Hintikka} labelling, while differing arbitrarily on unused template
indices, unused ports, unused fresh indices, or row-function arguments.  The
formalization has no finite normal form or quotient identifying these as the
same certificate.

\paragraph{Required repair.}
Introduce finite encoded syntax and make total evaluators derived from it.  A
repair can take either of these forms:
\begin{enumerate}
  \item Replace functional fields by bounded finite tables, for example
\begin{lstlisting}
structure FinTemplate where
  nSlots : Nat
  nFresh : Nat
  netTable : Array Atom       -- indexed by bounded ports

structure FinCatalog where
  nCore : Nat
  coreNetTable : Array Atom   -- indexed by Fin nCore x Fin nCore
  templates : Array FinTemplate
  root : Fin templates.size
  attaches : Array FinAttach
  FTable : Array Atom         -- indexed by template, row-code, row-code, fresh index
\end{lstlisting}
  and define \texttt{evalTemplate}, \texttt{evalCatalog}, \texttt{evalF} as
  total functions from these finite codes.

  \item Or keep the existing functional structures only as semantic views, but
  add a separate finite-code predicate:
\begin{lstlisting}
structure EncodedByFiniteCode (K : Catalog) (code : FinCatalogCode) : Prop := ...
\end{lstlisting}
  and require \texttt{CompletenessObligation} to produce such a code, with a
  computable size bound.
\end{enumerate}
In both approaches, \texttt{F\_reads} should become a theorem about a finite
row-code evaluator, not an unrestricted higher-order assumption.

\subsection{F2. \texttt{Satisfiable} fixes the carrier to \texttt{Occ}}

\textbf{Locations.}
\begin{itemize}
  \item Reference semantics: \texttt{ALCI\_RCC5\_REFERENCE.md}, lines 50--55
        and 70--73.
  \item \texttt{Interp}: lines 4129--4135.
  \item \texttt{Satisfiable}: lines 4280--4282.
\end{itemize}

The reference interpretation is \(I=(\Delta,\rho,\cdot^I)\), where
\(\Delta\) is an arbitrary nonempty domain of regions.  The Lean structure is:

\begin{lstlisting}
structure Interp where
  dom : Occ -> Prop
  rho : Occ -> Occ -> Atom
  val : Nat -> Occ -> Prop

def Satisfiable (C0 : Concept) : Prop :=
  exists I : Interp, RCC5Interp I /\ exists x, I.dom x /\ sat I x C0
\end{lstlisting}

Thus satisfiability is quantified only over subdomains of the countable carrier
\texttt{Occ}.  This is plausibly harmless for ordinary ALC-style concept
satisfiability if one proves a generated-submodel/countable-model property.  An
ALCI model can usually be restricted to the root and recursively chosen
existential witnesses; universal restrictions are downward-preserved, and RCC5
R1/R2/R3 are inherited by induced substructures.  But that theorem is not
stated here.

\paragraph{Required repair.}
Use a carrier-polymorphic interpretation:

\begin{lstlisting}
structure Interp (alpha : Type) where
  dom : alpha -> Prop
  rho : alpha -> alpha -> Atom
  val : Nat -> alpha -> Prop

def ReferenceSatisfiable (C0 : Concept) : Prop :=
  exists alpha, exists I : Interp alpha,
    RCC5Interp I /\ exists x, I.dom x /\ sat I x C0
\end{lstlisting}

Then prove that the certificate semantics embeds into this reference predicate.
Alternatively, keep \texttt{Occ} for the generated model but add an explicit
bridge theorem:

\begin{lstlisting}
theorem reference_satisfiable_iff_occ_satisfiable (C0 : Concept) :
  ReferenceSatisfiable C0 <-> Satisfiable C0 := ...
\end{lstlisting}

Until that bridge exists, the Lean target predicate is an under-approximation
of the reference semantics.

\subsection{F3. \texttt{CompletenessObligation} is not a finite decision statement}

\textbf{Location.} \texttt{CompletenessObligation}, lines 4459--4476.

The obligation says that every \texttt{Satisfiable} concept has some catalogue,
plan, list labelling, and root satisfying the generated soundness interface.
It does not assert any of the following:
\begin{itemize}
  \item a computable bound on the number of templates, core nodes, slots,
        fresh ports, plan length, or row codes;
  \item that catalogues and certificates are finite encoded data rather than
        arbitrary functions;
  \item that \texttt{tau} ranges only over the closure of \(C_0\);
  \item that \texttt{CatOk}, \texttt{PlanOk}, \texttt{SCat}, and
        \texttt{Hintikka} are decidable over the encoded search space;
  \item a theorem of the form \texttt{Decidable (ReferenceSatisfiable C0)} or
        an executable Boolean decision procedure with correctness proof.
\end{itemize}

The source comment at lines 4451--4458 mentions the finite type space
\(\leq 2^{|closure(C_0)|}\), but this bound is not part of the proposition.
As stated, \texttt{CompletenessObligation} is a semantic existence theorem over
higher-order Lean objects.  It can close a semantic equivalence, but not by
itself a finite search reduction.  In the project's terminology, this is best
understood as an underformalized bounded-width/F3 interface, not as an
additional independent completeness conjecture.

\paragraph{Required repair.}
Replace or supplement \texttt{CompletenessObligation} with a bounded finite
version.  A suitable shape is:

\begin{lstlisting}
def CompletenessBounded : Prop :=
  forall C0, ReferenceSatisfiable C0 ->
    exists code : FinCatalogCode C0,
      code.size <= bound C0 /\
      Decodable code /\
      let K := decodeCatalog code
      exists planCode tauCode rootCode,
        planCode.size <= planBound C0 /\
        tauCode.closedUnder (closure C0) /\
        CatOk K /\ PlanOk K (decodePlan planCode) /\
        SCat (buildCert K (decodePlan planCode)) /\
        Hintikka ... /\ rootCarries C0
\end{lstlisting}

Then add the actual decision theorem:

\begin{lstlisting}
theorem decidable_reference_satisfiable (C0 : Concept) :
  Decidable (ReferenceSatisfiable C0) := ...
\end{lstlisting}

or, preferably, an executable checker:

\begin{lstlisting}
def decideSat (C0 : Concept) : Bool := ...

theorem decideSat_correct (C0 : Concept) :
  decideSat C0 = true <-> ReferenceSatisfiable C0 := ...
\end{lstlisting}

This is the most important repair for the advertised decidability claim.

\subsection{F4. \texttt{SCat.net\_r3} over-checks degenerate template triples}

\textbf{Locations.}
\begin{itemize}
  \item \texttt{SCond.tri}: lines 1861--1866.
  \item \texttt{SCat.net\_r3}: lines 2797--2802.
\end{itemize}

The semantic S-condition excludes all degenerate triples through a fresh
occurrence:

\begin{lstlisting}
SCond.tri : ...
  x != y -> x != born i jz -> y != born i jz ->
  pairVal C x y in comp (pairVal C x (born i jz))
                       (pairVal C (born i jz) y)
\end{lstlisting}

The catalogue check has only \texttt{p != q}:

\begin{lstlisting}
SCat.net_r3 : ...
  p in memberPortsList C i -> q in memberPortsList C i ->
  j < nFresh -> p != q ->
  net p q in comp (net p fresh_j) (conv (net q fresh_j))
\end{lstlisting}

It does not require \texttt{p != fresh\_j} or \texttt{q != fresh\_j}.  Hence
it can inspect \texttt{net fresh\_j fresh\_j}, a diagonal template value that
is not used as an off-diagonal semantic relation.  This makes \texttt{SCat}
strictly stronger than the intended local R3 check unless template diagonals are
normalized.

\paragraph{Concrete calculation.}
Consider one template with two fresh ports and no core/slots.  Set:
\begin{lstlisting}
net fresh0 fresh1 = PP
net fresh1 fresh0 = PPI
net fresh0 fresh0 = DR
net fresh1 fresh1 = DR
\end{lstlisting}
Converse coherence holds for these entries.  Semantically, \texttt{SCond.tri}
has no nondegenerate triple among only two fresh occurrences.  But
\texttt{SCat.net\_r3}, with \texttt{p = fresh0}, \texttt{q = fresh1}, and
\texttt{j = 1}, demands:
\begin{lstlisting}
PP in comp PP (conv DR)
\end{lstlisting}
Since \texttt{conv DR = DR} and \texttt{comp PP DR = [DR]}, this is false.
Thus \texttt{SCat} rejects an object for an irrelevant diagonal choice.

\paragraph{Required repair.}
Use one of the following fixes.
\begin{enumerate}
  \item Align \texttt{SCat.net\_r3} with \texttt{SCond.tri} by adding side
        conditions:
\begin{lstlisting}
p != inr (inr j) -> q != inr (inr j) ->
net p q in comp (net p fresh_j) (conv (net q fresh_j))
\end{lstlisting}
  \item Alternatively, require normalized template diagonals:
\begin{lstlisting}
net_diag : forall t p, validPort t p -> (K.tmpl t).net p p = EQ
\end{lstlisting}
        and include this in \texttt{CatOk}/\texttt{Wellformed}.  Then make
        explicit that templates are total atomic networks including diagonal
        EQ, not just off-diagonal ledgers.
\end{enumerate}
The first repair is semantically minimal.  The second is also acceptable if
normalization is intended.

\subsection{F5. Current positive witnesses are not rich enough}

\textbf{Locations.}
\begin{itemize}
  \item \texttt{certC}: lines 1657--1664.
  \item \texttt{certK}: lines 3995--4003.
  \item \texttt{sample\_satisfiable\_ex}: lines 4483--4534.
\end{itemize}

The provided positive certificate/catalogue witnesses use \texttt{nCore := 0}
and all-DR template/core/steering data.  They exercise plumbing and inherited
slot paths but not the hard positive cases: PP/PPI transitivity, PO cases,
core rows, nontrivial steering, and non-DR successful composition.
\texttt{sample\_satisfiable\_ex} does exercise a real existential, but it is a
hand-built two-point interpretation, not a generated certificate plus Hintikka
labelling passed through \texttt{sat\_from\_hintikka}.

This is not evidence of vacuity.  A simple rich positive certificate should be
available: three fresh occurrences with
\[
  0\;PP\;1, \qquad 1\;PP\;2, \qquad 0\;PP\;2,
\]
converses in the opposite direction and diagonal EQ.  This exercises
\texttt{comp PP PP = [PP]} positively.

\paragraph{Required repair.}
Add at least one generated positive witness with:
\begin{itemize}
  \item three fresh points and a PP-chain using \texttt{comp PP PP};
  \item at least one PO case;
  \item at least one inherited slot;
  \item nonzero core or an explicit explanation why core-free certificates
        suffice for the witness suite;
  \item a root-anchored \texttt{Hintikka} labelling discharged through
        \texttt{sat\_from\_hintikka} or \texttt{generated\_satisfiable}.
\end{itemize}

\section{Adequacy checks with no defect found}

\subsection{Syntax}

\textbf{Location.} \texttt{Concept}, lines 4117--4127.

The constructors match the reference grammar in NNF:
\begin{lstlisting}
top | bot | atom | natom | and | or | ex Atom | all Atom
\end{lstlisting}
There is no general negation constructor, which correctly builds the NNF
discipline into the syntax.  Roles are exactly RCC5 atoms; inverse roles are
absorbed by converse, so \(PP^{-}=PPI\) is representable without a separate
inverse-role constructor.

\subsection{Satisfaction}

\textbf{Location.} \texttt{sat}, lines 4137--4146.

The role reading and quantifier polarity are correct:
\begin{lstlisting}
| x, ex r c  => exists y, I.dom y /\ I.rho x y = r /\ sat I y c
| x, all r c => forall y, I.dom y -> I.rho x y = r -> sat I y c
\end{lstlisting}
This matches the reference reading \(R(x,y)\) iff \(\rho(x,y)=R\).  The guard
is on the quantified element \(y\).  The only caveat is that \texttt{sat I x C}
is defined for any \texttt{x : Occ}; theorem statements that use it as model
truth carry a domain hypothesis, so this is not a soundness bug.

\subsection{RCC5 interpretation frame}

\textbf{Location.} \texttt{RCC5Interp}, lines 4206--4214.

The fields encode reflexive EQ, strong EQ as identity, converse coherence, and
R3 composition closure:
\begin{lstlisting}
refl_eq : dom x -> rho x x = EQ
eq_id   : dom x -> dom y -> rho x y = EQ -> x = y
conv_   : dom x -> dom y -> rho y x = conv (rho x y)
comp_   : dom x -> dom y -> dom z ->
          rho x z in comp (rho x y) (rho y z)
\end{lstlisting}
The R3 clause is stated for all triples, not just distinct triples, but the EQ
identity rows make the degenerate cases harmless in actual interpretations.
Standalone \texttt{RCC5Interp} permits an empty domain; \texttt{Satisfiable}
supplies the root witness.

\subsection{RCC5 table}

\textbf{Location.} \texttt{conv} and \texttt{comp}, lines 138--159.

The converse table and the relevant composition entries agree with the standard
RCC5 atomic composition table, including:
\begin{lstlisting}
comp PP PP   = [PP]
comp PP DR   = [DR]
comp DR PPI  = [DR]
comp PPI PP  = [EQ, PP, PPI, PO]
comp PO PO   = all atoms
comp DR DR   = all atoms
\end{lstlisting}
I found no table-orientation defect.

\section{Headline theorem statement audit}

\subsection{\texttt{frame\_closed}}

\textbf{Location.} Lines 2291--2300.

The theorem is conditional on operational frame entries existing:
\begin{lstlisting}
Frame.get? (unfoldAll C) x y = some v1
Frame.get? (unfoldAll C) x z = some v2
Frame.get? (unfoldAll C) z y = some v3
\end{lstlisting}
Therefore it is not by itself a theorem that the unfolded frame is a total
atomic network.  This is acceptable because later model construction uses
\texttt{pairVal}; the prose should not overstate \texttt{frame\_closed} as a
total-frame closure theorem.

\subsection{\texttt{scat\_scond}}

\textbf{Location.} Lines 3205--3227.

The statement is good:
\begin{lstlisting}
Wellformed C -> SCat C -> Faithful C -> SCond C
\end{lstlisting}
It does not smuggle in \texttt{SCond}.  The caveat is the over-strong
\texttt{SCat.net\_r3} clause described in F4.

\subsection{\texttt{catalogue\_soundness}}

\textbf{Location.} Lines 3946--3958.

The theorem is sound but has extra hypotheses beyond a slogan like
``\texttt{SCat} alone yields \texttt{SCond}'':
\begin{lstlisting}
CatOk K -> hcc -> hcr -> PlanOk K plan ->
SCat (buildCert K plan) -> SCond (buildCert K plan)
\end{lstlisting}
The converse-coherence hypothesis \texttt{hcc} and core-row hypothesis
\texttt{hcr} should either be folded into \texttt{CatOk} or named as part of
the generated-catalogue wellformedness predicate.

\subsection{\texttt{build\_wellformed} and \texttt{build\_faithful}}

\textbf{Locations.} \texttt{build\_wellformed}, line 3708; \texttt{build\_faithful}, line 3815.

\texttt{build\_wellformed} has the expected form:
\begin{lstlisting}
CatOk K -> PlanOk K plan -> Wellformed (buildCert K plan)
\end{lstlisting}
\texttt{build\_faithful} also requires global \texttt{hcc}/\texttt{hcr}
hypotheses.  This is not unsound, but those hypotheses are part of the finite
catalogue data and should be included in the eventual encoded check.

\subsection{\texttt{truth\_lemma} and \texttt{sat\_from\_hintikka}}

\textbf{Locations.} \texttt{truth\_lemma}, lines 4173--4197;
\texttt{sat\_from\_hintikka}, lines 4292--4300.

These statements are clean.  \texttt{sat\_from\_hintikka} is the central
soundness theorem:
\begin{lstlisting}
Wellformed C -> SCond C -> Hintikka dom rho tau ->
root in domain -> C0 in tau root -> Satisfiable C0
\end{lstlisting}
It neither assumes the conclusion nor hides the root witness.

\subsection{\texttt{model\_hintikkaP} and \texttt{satisfiable\_iff\_hintikkaP}}

\textbf{Locations.} Lines 4383--4405.

These are correct but weak relative to the prose phrase ``type-system
certificate loses nothing''.  The forward direction of
\texttt{satisfiable\_iff\_hintikkaP} chooses:
\begin{lstlisting}
tau x C := sat I x C
\end{lstlisting}
So the theorem proves unrestricted semantic Hintikka-realizability, not finite
closure-bounded type realizability, not generated-certificate realizability, and
not a decision procedure.  The theorem is useful as a semantic abstraction
lemma, but it should not be advertised as finite type-system completeness.

\subsection{\texttt{generated\_satisfiable}}

\textbf{Location.} Lines 4428--4449.

This is a sound composition theorem: a generated catalogue/plan satisfying the
listed conditions and carrying a root-anchored \texttt{Hintikka} labelling
produces \texttt{Satisfiable C0}.  It is not itself a completeness theorem; it
is the target interface for the missing generator.

\section{Required repairs as a consolidated checklist}

\begin{enumerate}
  \item \textbf{Separate reference semantics from generated semantics.}
        Define carrier-polymorphic \texttt{ReferenceSatisfiable}, or prove an
        explicit countable/\texttt{Occ}-model theorem bridging arbitrary
        reference models to the generated carrier.

  \item \textbf{Finite-code the certificate and catalogue syntax.}
        Replace arbitrary \texttt{Nat -> Nat -> Atom}, \texttt{TPort -> TPort -> Atom},
        and higher-order \texttt{F} fields by finite arrays or finite codes,
        with total evaluators and proofs that all semantic views are decoded
        from those codes.

  \item \textbf{Strengthen \texttt{CompletenessObligation}.}
        State a bounded obligation that produces finite codes of size bounded
        by computable functions of \(C_0\), closure-bounded type labels, and
        finite plan/root codes.  This is the place where the bounded-width/F3
        mathematical assumption should be formalized.  Add a proved checker or
        \texttt{Decidable}-theorem for the bounded search space.

  \item \textbf{Fix \texttt{SCat.net\_r3}.}
        Add \texttt{p != fresh\_j} and \texttt{q != fresh\_j}, or require
        normalized diagonal EQ for all valid template ports and document that
        templates are total atomic networks with diagonal EQ.

  \item \textbf{Fold global catalogue coherence into \texttt{CatOk}.}
        The hypotheses \texttt{hcc} and \texttt{hcr} occur in
        \texttt{build\_faithful}, \texttt{catalogue\_soundness},
        \texttt{generated\_satisfiable}, and \texttt{CompletenessObligation}.
        They should be part of a single decidable finite catalogue-wellformedness
        predicate.

  \item \textbf{Add rich non-vacuity witnesses.}
        Include generated positive certificates exercising non-DR composition,
        PP/PPI chains, PO, inherited slots, and a real Hintikka existential
        discharged through the certificate pipeline.

  \item \textbf{Revise prose claims and theorem comments.}
        Keep the strong claim for the soundness pipeline.  Qualify
        \texttt{satisfiable\_iff\_hintikkaP} as semantic Hintikka completeness,
        not finite type-system completeness.  State explicitly that the
        current \texttt{CompletenessObligation} is not yet the finite bounded
        obligation needed for decidability.
\end{enumerate}

\section{What survives the review}

The important positive result is that the formalized soundness theorem looks
semantically meaningful.  A wellformed, S-conditioned certificate with a
Hintikka labelling yields an RCC5 interpretation and a root satisfying the
input concept.  The syntax and satisfaction clauses line up with the reference
logic.  The RCC5 frame predicate is the right strong-EQ atomic-frame notion.

Thus the defects do not suggest that the kernel has checked a vacuous proof of
an obviously wrong satisfaction relation.  They instead identify the missing
finite and reference-semantics bridges needed before the artifact can honestly
claim a finite decidability reduction.

\appendix
\section{Accompanying scripts}

The package contains the following scripts.
\begin{itemize}
  \item \texttt{scripts/audit\_static.py}: extracts source locations and flags
        the higher-order fields and theorem signatures discussed above.
  \item \texttt{scripts/scat\_degenerate\_check.py}: computes the concrete RCC5
        table calculation behind F4.
  \item \texttt{scripts/build\_report.sh}: recompiles this PDF from the TeX
        source.
  \item \texttt{scripts/optional\_lean\_check.sh}: optionally invokes Lean on
        \texttt{Round19Transport.lean} when Lean is installed.
\end{itemize}

\end{document}
